% Part: many-valued-logics
% Chapter: sequent-calculus
% Section: proof-theoretic-notions

\documentclass[../../../include/open-logic-section]{subfiles}

\begin{document}

\olfileid{mvl}{seq}{prf}

\olsection{证明论概念}

\begin{explain}
正如我们已经定义了若干重要的语义概念（有效性、语义后承、可满足性）一样，我们现在来定义对应的\emph{证明论概念}。这些概念不是通过语句在结构中的满足来定义的，而是借助某些矢列的可推导性或不可推导性来定义的。这些概念是重合的，这是一个重要的发现。这种重合正是\emph{可靠性定理与完备性定理}的内容。
\end{explain}

\begin{defn}[定理]
若在~$\Log{LK}$ 中存在矢列 $\quad \Sequent !A$ 的推导，则称语句~$!A$ 为\emph{定理}。若 $!A$ 是定理，则记作 $\Proves !A$；若不是，则记作 $\Proves/ !A$。
\end{defn}

\begin{defn}[可推导性]
一个语句~$!A$ \emph{可以从}语句集~$\Gamma$ \emph{推导}，记作 $\Gamma \Proves !A$，当且仅当存在一个有限子集~$\Gamma_0 \subseteq \Gamma$，以及一个由~$\Gamma_0$ 中语句构成的序列 $\Gamma_0'$，使得 $\Log{LK}$ 能推导出 $\Gamma_0' \Sequent !A$。若 $!A$ 不能从 $\Gamma$ 推导，则记作 $\Gamma \Proves/ !A$。
\end{defn}

由于收缩、弱化和交换规则的存在，$\Gamma_0'$ 中语句的顺序和数量无关紧要：如果矢列 $\Gamma_0' \Sequent !A$ 是可推导的，那么对任何由与~$\Gamma_0'$ 相同语句构成的 $\Gamma_0''$，矢列 $\Gamma_0'' \Sequent !A$ 也是可推导的。例如，若 $\Gamma_0 = \{!B, !C\}$，
那么 $\Gamma_0' = \tuple{!B, !B, !C}$ 和 $\Gamma_0'' =
\tuple{!C, !C, !B}$ 都是仅包含~$\Gamma_0$ 中语句的序列。如果以其中一个序列为前件的矢列是可推导的，那么另一个也可推导，例如：
\begin{prooftree}
  \AxiomC{}
  \Deduce$!B, !B, !C \fCenter !A$
  \RightLabel{\LeftR{\Contraction}}
  \UnaryInf$!B, !C \fCenter !A$
  \RightLabel{\LeftR{\Exchange}}
  \UnaryInf$!C, !B \fCenter !A$
  \RightLabel{\LeftR{\Weakening}}
  \UnaryInf$!C, !C, !B \fCenter !A$
\end{prooftree}

从现在起，如果 $\Gamma_0$ 是一个有限语句集，我们将用 $\Gamma_0 \Sequent !A$ 表示任一前件为~$\Gamma_0$ 中语句序列的矢列，并在必要时默认允许使用收缩、交换和弱化规则。

\begin{defn}[一致性]
一个语句集~$\Gamma$ 是\emph{不一致的}，当且仅当存在一个有限子集~$\Gamma_0 \subseteq \Gamma$，使得 $\Log{LK}$ 能推导出 $\Gamma_0 \Sequent \quad$。若 $\Gamma$ 不是不一致的，即对每个有限子集 $\Gamma_0 \subseteq \Gamma$，$\Log{LK}$ 都不能推导出 $\Gamma_0 \Sequent \quad$，则称它是\emph{一致的}。
\end{defn}

\begin{prop}[自反性]
\ollabel{prop:reflexivity}
若 $!A \in \Gamma$，则 $\Gamma \Proves !A$。
\end{prop}

\begin{proof}
初始矢列 $!A \Sequent !A$ 是可推导的，且 $\{!A\} \subseteq \Gamma$。
\end{proof}

\begin{prop}[单调性]
\ollabel{prop:monotonicity}
若 $\Gamma \subseteq \Delta$ 且 $\Gamma \Proves !A$，则 $\Delta \Proves !A$。
\end{prop}

\begin{proof}
假设 $\Gamma \Proves !A$，即存在有限子集 $\Gamma_0 \subseteq \Gamma$，使得 $\Gamma_0 \Sequent !A$ 可推导。由于 $\Gamma \subseteq \Delta$，则 $\Gamma_0$ 也是~$\Delta$ 的一个有限子集。因此，$\Gamma_0 \Sequent !A$ 的推导也表明 $\Delta \Proves !A$。
\end{proof}

\begin{prop}[传递性]
\ollabel{prop:transitivity}
若 $\Gamma \Proves !A$ 且 $\{!A\} \cup \Delta \Proves
!B$，则 $\Gamma \cup \Delta \Proves !B$。
\end{prop}

\begin{proof}
如果 $\Gamma \Proves !A$，则存在有限子集 $\Gamma_0 \subseteq \Gamma$ 以及 $\Gamma_0 \Sequent !A$ 的一个推导~$\pi_0$。如果 $\{!A\}
\cup \Delta \Proves !B$，则对于某个有限子集 $\Delta_0 \subseteq \Delta$，存在一个 $!A, \Delta_0 \Sequent !B$ 的推导~$\pi_1$。考虑如下推导：
\begin{prooftree}
  \AxiomC{}
  \RightLabel{$\pi_0$}
  \Deduce$\Gamma_0 \fCenter !A$
  \AxiomC{}
  \RightLabel{$\pi_1$}
  \Deduce$!A, \Delta_0 \fCenter !B$
  \RightLabel{\Cut}
  \BinaryInf$\Gamma_0, \Delta_0 \fCenter !B$
\end{prooftree}
由于 $\Gamma_0 \cup \Delta_0 \subseteq \Gamma \cup \Delta$，这便证明了 $\Gamma \cup \Delta \Proves !B$。
\end{proof}

注意，这特别意味着，若 $\Gamma \Proves !A$ 且 $!A \Proves !B$，则 $\Gamma \Proves !B$。同样可得，若 $!A_1,
\dots, !A_n \Proves !B$ 且对每个~$i$ 有 $\Gamma \Proves !A_i$，则 $\Gamma \Proves !B$。

\begin{prop}
\ollabel{prop:incons}
$\Gamma$ 是不一致的，当且仅当 $\Gamma \Proves {!A}$ 对每个语句~$!A$ 成立。
\end{prop}

\begin{proof}
习题。
\end{proof}

\begin{prob}
证明 \olref[fol][seq][ptn]{prop:incons}
\end{prob}

\begin{prop}[紧致性]
  \ollabel{prop:proves-compact}
  \begin{enumerate}
  \item 若 $\Gamma \Proves !A$，则存在有限子集 $\Gamma_0 \subseteq \Gamma$，使得 $\Gamma_0 \Proves !A$。
  \item 若~$\Gamma$ 的每个有限子集都是一致的，则 $\Gamma$ 是一致的。
  \end{enumerate}
\end{prop}

\begin{proof}
  \begin{enumerate}
    \item 若 $\Gamma \Proves !A$，则存在有限子集 $\Gamma_0 \subseteq \Gamma$，使得矢列 $\Gamma_0 \Sequent !A$ 有一个推导。因此，$\Gamma_0 \Proves !A$。
    \item 若 $\Gamma$ 不一致，则存在有限子集 $\Gamma_0 \subseteq \Gamma$，使得 $\Log{LK}$ 能推导出 $\Gamma_0 \Sequent \quad$。但这样一来，$\Gamma_0$ 就是 $\Gamma$ 的一个不一致的有限子集。
  \end{enumerate}
\end{proof}

\end{document}